\section{Local markers and episode ends}\label{sec:geometry}

The split identity requires each accepted suffix to stop at the actual
first episode end. Agreement of tags alone cannot guarantee that endpoint.
This section constructs an intrinsic episode rule from local markers and
proves two ways for a suffix to identify its end, including when its guessed
origin is false. The same marker rule also identifies periodic intervals
for the second outer split.

We use integer positions between letters: the letter at position $x$
occupies $[x,x+1)$.  Intervals of possible marker positions, in contrast,
are inclusive integer intervals.  Throughout this section the alphabet is
nonempty, so every finite word has a two-sided extension.  Such extensions
will only be used to prove statements about local rules; all tests on a
finite word must have their stated inspection intervals inside that word.

\subsection{An elementary local marker rule}

A positive integer $d$ is a \emph{period} of a finite block if letters at
positions $d$ apart agree whenever both positions belong to the block.
Equivalently, the block extends to a two-sided $d$-periodic word.  If the
block has length at least $d$, this extension is unique for that period.

This use of markers to isolate local periodicity follows the viewpoint of
Krieger's marker lemma~\cite[Lemma~2]{Krieger1982}. Related clopen-marker
constructions are set out in \cite[Lemmas~3.2--3.4]{Meyerovitch2025}.
The local rule and its explicit finite-window bounds are proved below.

\begin{lemma}[Local markers]\label{lem:markers}
For integers $d\ge2$ and $K\ge1$ there are an integer $r\ge K$ and a fixed,
translation-equivariant marking rule on two-sided words with the following
properties.
\begin{enumerate}[label=\textnormal{(\roman*)}]
\item Distinct markers have distance at least $d$.
\item Whether $x$ is marked depends only on the letters in $[x-r,x+r)$.
\item If no marker lies in $[z-d,z+d]$, the block $[z,z+K)$ has a period
strictly less than $d$.
\end{enumerate}
Consequently a marker test whose inspection interval is present in a
finite word is independent of its two-sided extension.
\end{lemma}

\begin{proof}
List as $v_1,\ldots,v_N$ all length-$K$ words having no period below $d$.
Start with no markers.  At stage $j$, mark every occurrence of $v_j$ whose
start is at distance at least $d$ from all markers installed at earlier
stages.  Two occurrences installed at this same stage cannot be less than
$d$ apart.  Indeed a word in the list has $K\ge d$, since otherwise $K$
itself would be a period below $d$.  Starts at positive distance
$a<d\le K$ would therefore overlap and give $v_j$ period $a$, a
contradiction.  Separation is preserved at every stage.

For locality, take $r_0=K$ and $r_j=K+j(d-1)$.  If earlier decisions have
radius $r_{j-1}$, deciding the stage-$j$ rule at $x$ requires its length-$K$
block and the earlier decisions at $x-(d-1),\ldots,x+(d-1)$.  These are all
determined in $[x-r_j,x+r_j)$.  The union with earlier markers is determined
there too.  Thus $r=r_N$ suffices.  The construction commutes with every
integer translation.

If the block at $z$ equals some $v_j$, at stage $j$ its start is either
marked or blocked by an earlier marker at distance less than $d$.
That marker remains installed.  Absence of a marker in $[z-d,z+d]$
therefore excludes every word on the list and proves the last assertion.
Locality proves extension independence.
\end{proof}

\begin{lemma}[Overlap of short-period blocks]\label{lem:periodic-overlap}
Two two-sided words of respective periods $a,b$ which agree on $ab$
consecutive letter positions agree everywhere.  In particular:
\begin{enumerate}[label=\textnormal{(\roman*)}]
\item If $K-w>d^2$, length-$K$ blocks starting at a finite nonempty set
of positions of diameter at most $w$, each having a period below $d$,
have one common two-sided periodic extension.  The word agrees with this
extension on the whole union of these blocks.
\item Use the rule of Lemma~\ref{lem:markers} with $K>3d^2$.
If $H\ge d$ and $[z-H,z+H]$ contains no marker, all the blocks starting
at integer positions from $z-H+d$ through $z+H-d$ belong to one common
two-sided periodic word of some period below $d$.  The original word
agrees with it throughout $[z-H+d,z+H-d+K)$.
\end{enumerate}
All statements apply to a finite word when the indicated blocks and
marker inspection intervals are available.
\end{lemma}

\begin{proof}
The integer $ab$ is a period of both two-sided words.  Given any position,
translate it by a multiple of $ab$ into the interval of agreement.  Both
letters are preserved by this translation, proving the first assertion.

For (i), any two blocks overlap in at least $K-w>d^2$ letters, more than
the product of their two short periods.  Their periodic extensions
therefore agree.  This also shows that choosing a different short period
for any one block gives the same extension.  Fixing one block now gives
a common extension for them all.  Their union is the interval from the
least start to the greatest start plus $K$, since $w<K$; every letter
there belongs to a block and has the asserted value.

For (ii), each indicated start $x$ has
$[x-d,x+d]\subseteq[z-H,z+H]$, so Lemma~\ref{lem:markers} gives a period
below $d$ for its length-$K$ block.  Consecutive blocks overlap in $K-1$
letters, at least the product of their periods.  The first assertion,
applied successively, identifies all their extensions.  The common
extension retains, for example, the short period of the first block.
\end{proof}

\subsection{How often a sliding search can fail}

For a set $S\subseteq\mathbb Z$ of marker positions and an integer
$L\ge0$, write
\[
 f_L(z)=\min\bigl(S\cap[z-L,z+L]\bigr),
\]
with value $\bot$ when the intersection is empty.  Equality of two search
results means equality of their positions in the word, or that both are
$\bot$.

\begin{lemma}[Search crossings and right margins]\label{lem:search-crossings}
Let $I$ be an interval of integer centers of diameter $D$.
\begin{enumerate}[label=\textnormal{(\roman*)}]
\item Suppose distinct points of $S$ are at least $d$ apart.  If $h\ge0$
and $D+2h<d$, at most $2(2h+1)$ centers $z\in I$ satisfy
\[
 f_L(z+e)\ne f_L(z)\quad\text{for some integer }|e|\le h.
\]
This estimate is independent of $L$.
\item Let $h\ge1$.  If $D<2L+2$, at most $h$ centers $z\in I$ satisfy
\[
 S\cap[z-L,z+L]=\varnothing,
 \qquad S\cap(z+L,z+L+h]\ne\varnothing.
\]
This estimate does not require a separation hypothesis on $S$.
\end{enumerate}
If centers are indexed by ordered pairs $(p,p')$, $p\ne p'$, as
$z=z_0+p-p'$, and all their ordered nonzero differences are distinct,
the same bounds count ordered pairs.
\end{lemma}

\begin{proof}
If a first-marker result changes, a marker belongs to the symmetric
difference of the two search intervals.  Such a marker is within distance
$h$ of $z-L$ or of $z+L$.  As $z$ ranges over $I$, each of these enlarged
endpoint ranges has diameter $D+2h<d$ and hence contains at most one
marker.  For a fixed marker and a fixed endpoint, the condition of being
within distance $h$ permits at most $2h+1$ integer centers.  Summing over
the two endpoints proves (i).  We only use the implication from a changed
result to a crossing; an entering marker need not change the minimum.

For (ii), assign to a center $z$ its first marker $m$ strictly after
$z+L$.  It satisfies $z+L<m\le z+L+h$.  Suppose distinct serving markers
$m<m'$ are assigned to centers $z,z'$.  Since $m'$ is the first marker
after $z'+L$, we have $m\le z'+L$.  The marker-free interval about $z'$
then forces $m\le z'-L-1$.  Combining this with $m\ge z+L+1$ gives
$z'-z\ge2L+2$, contrary to the diameter hypothesis.  Thus at most one
marker serves the centers under consideration.  For that fixed $m$,
\[
 m-L-h\le z\le m-L-1,
\]
which allows at most $h$ integer centers.  Finally, distinct ordered
differences make the indexing map $(p,p')\mapsto z_0+p-p'$ injective.
\end{proof}

For local markers, the counting estimates apply to a finite word when
all the shifted search inspection intervals are present.  Indeed one
may evaluate the rule on any common two-sided extension, and locality
identifies all the results being counted.

The marker and search lemmas now supply the two kinds of control we need:
absence yields a periodic interval, while instability can occur at only a
bounded number of translated centers. We next use a separate marker search
to define an episode end that an independently chosen suffix can check.

\subsection{The episode rule}

We now introduce the geometric parameters; their finite order of choice
and all placement requirements are verified in
Proposition~\ref{prop:parameters}.  An episode begins at $s$, and its
sampling origin is
\[
 t=s+L_0.
\]
The five aligned subbuffers will lie between $s$ and $t$.  The outer
split layers will use finitely many nominal offsets after $t$.

Apply Lemma~\ref{lem:markers} to an \emph{inner} marker system with
parameters
\[
 d_\circ\ge2,\qquad K_\circ>3d_\circ^2,
 \qquad r_\circ\ge K_\circ.
\]
For each label $\lambda$ in a finite set $\mathcal Z$, let $o_\lambda$
be its nominal offset and set
\[
 z_\lambda(t)=t+o_\lambda,\qquad
 u_\lambda(t)=\text{first inner marker in }
 [z_\lambda(t)-H,z_\lambda(t)+H],
\]
again using $\bot$ for absence.  The labels, offsets, trials, and full
cut-offset sets $\mathcal H_1,\mathcal H_2$ are specified in
Section~\ref{sec:outer}; no trial roster is needed to define the end rule.

Independently apply Lemma~\ref{lem:markers} to an \emph{end} marker system
with parameters
\[
 d_\bullet\ge2,\qquad K_\bullet>3d_\bullet^2,
 \qquad r_\bullet\ge K_\bullet.
\]
Its nominal position and search are
\[
 z_\bullet(t)=t+o_\bullet,\qquad
 \eta(t)=\text{first end marker in }
 [z_\bullet(t)-d_\bullet,z_\bullet(t)+d_\bullet].
\]
The value is $\bot$ in the markerless case.  The inspection interval for
this search is
\[
 [z_\bullet(t)-d_\bullet-r_\bullet,
   z_\bullet(t)+d_\bullet+r_\bullet).
\]
When $\eta(t)=\bot$, Lemma~\ref{lem:markers} gives the seed block
$[z_\bullet(t),z_\bullet(t)+K_\bullet)$ a period below $d_\bullet$.
Lemma~\ref{lem:periodic-overlap} makes its two-sided periodic extension
unique even when several such periods are possible.  Denote it by
$\pi_t$.

\begin{definition}[Complete episodes]\label{def:episodes}
Fix a positive integer $B_{\rm tail}$.  With all required bounded data
present, the \emph{anchor} at origin $t$ is defined as follows:
\begin{enumerate}[label=\textnormal{(\roman*)}]
\item if $\eta(t)\ne\bot$, the anchor is $q=\eta(t)$;
\item if $\eta(t)=\bot$, the anchor is the first letter position
$q\ge z_\bullet(t)+K_\bullet$ whose letter differs from $\pi_t(q)$.
If no such letter is present, no anchor has yet been found.
\end{enumerate}
An episode beginning at $s$ ends exactly at $q+B_{\rm tail}$.
The language $E$ consists of the nonempty words which, with $s=0$, contain
the required search and seed data, have an anchor under this rule, and
end exactly there.  In the markerless case the tail includes the mismatch
letter at $q$, since $B_{\rm tail}\ge1$.
\end{definition}

In the markerless case, ``first'' is part of the end test: every letter
from the seed's end up to $q$ must agree with the repetition, and the
letter at $q$ must disagree.  Applying the rule at an origin relative to a
suffix beginning at a cut uses precisely the same convention.

\begin{lemma}[Prefix code]\label{lem:prefix-code}
The language $E$ is a prefix code.  Its end decisions have the same
values in an episode considered alone and in any word containing that
episode at the same position.
\end{lemma}

\begin{proof}
Let $e,f\in E$, with $e$ a prefix of $f$.  Their starts and sampling
origins coincide.  All end-marker inspections used by $e$ are internal
to $e$, so the two words obtain the same value of $\eta(t)$.  A present
marker immediately fixes the same anchor.  If it is absent, the same
seed fixes the same repetition.  The first mismatch seen in $e$ is
present in $f$, and all earlier letters inspected for continuation
coincide; hence their first mismatches coincide too.  In either case
$|e|=q+B_{\rm tail}=|f|$, proving the prefix-code assertion.  The same
argument, and locality of all bounded searches, prove the final claim.
\end{proof}

\subsection{Independent suffix choices have one end}

Fix a stability radius satisfying
\begin{equation}\label{eq:end-overlap}
 0\le w_E<d_\bullet,
 \qquad K_\bullet-w_E>d_\bullet^2.
\end{equation}
We say that $\eta$ is \emph{stable within $w_E$ at $t$} if
$\eta(t+v)=\eta(t)$ for every integer $|v|\le w_E$.

\begin{lemma}[Shared anchor and exact suffix consumption]
\label{lem:shared-anchor}
Let a word begin with a complete episode $e$ at $s$, with origin $t$,
anchor $q$, and end $b_0=q+B_{\rm tail}$.  Let $s<k<b_0$ be a cut, and
let $[k,b)$ be any proposed suffix factor of that word; its endpoint $b$
may be before or after $b_0$.  Let $\mathcal T$ be a finite nonempty set
of candidate origins such that
\[
 t\in\mathcal T,\qquad \operatorname{diam}\mathcal T\le w_E.
\]
Assume that every candidate's end-search inspections and seed block lie
inside both $[s,b_0)$ and $[k,b)$.  Suppose the proposed suffix requires
\begin{enumerate}[label=\textnormal{(\roman*)}]
\item equality of $\eta(\tau)$ for all $\tau\in\mathcal T$;
\item the exact end rule of Definition~\ref{def:episodes} at any one
reference $\tau_0\in\mathcal T$, with anchor $b-B_{\rm tail}$.
\end{enumerate}
Then $b=b_0$, and the reference anchor is $q$.
Moreover, if $\eta$ is stable within $w_E$ at $t$ and all the candidate
data are visible in the actual suffix $[k,b_0)$, these two requirements
hold for that suffix and every choice of reference.
\end{lemma}

\begin{proof}
All bounded tests in the two factors inspect the same ambient letters.
By locality they obtain the same results, including for false candidate
origins.  If the common value of $\eta$ is a marker position, it is also
$\eta(t)=q$, so the exact end requirement gives $b=q+B_{\rm tail}$.

Suppose instead that every value is $\bot$.  The seed starts
$x_\tau=z_\bullet(\tau)$ have diameter at most $w_E$.
By Equation~\eqref{eq:end-overlap} and
Lemma~\ref{lem:periodic-overlap}, all seeds have one common two-sided
extension $\pi$.  The word agrees with $\pi$ throughout their union
\[
 [x_{\min},x_{\max}+K_\bullet),
 \qquad
 x_{\min}=\min_{\tau\in\mathcal T}x_\tau,
 \quad x_{\max}=\max_{\tau\in\mathcal T}x_\tau.
\]
Thus the true first mismatch $q$, which occurs after the true seed,
cannot lie inside any later seed.  It follows that
$q\ge x_{\max}+K_\bullet$.  There is no mismatch between the end of any
seed and $q$: before $x_{\max}+K_\bullet$ this follows from the seed
union, and afterwards it follows from the true first-mismatch rule.
Hence every candidate seed has exactly the same first subsequent
mismatch $q$.

For completeness this argument also excludes incorrectly truncated or
extended proposed suffixes.  Put $q'=b-B_{\rm tail}$.  If $q'<q$, the
proposed mismatch letter at $q'$ is present in both factors and
contradicts the continuation just proved.  If $q'>q$, the proposed
continuation reaches the visible mismatch at $q$ and fails there.  A
suffix lacking any required bounded inspection or its claimed mismatch
letter cannot pass the requirements in the first place.  Therefore
$q'=q$ and $b=b_0$.

Finally, stability and $t\in\mathcal T$ imply equality of $\eta$ on
$\mathcal T$.  The same seed-union argument shows that the actual suffix
satisfies the exact end rule for every reference.
\end{proof}

Figure~\ref{fig:common-seeds} shows the markerless part of this argument:
all seed blocks lie in one common periodic word, and the first mismatch
occurs after their union. This is what makes the endpoint independent of
the chosen reference seed.

\begin{figure}[t]
\centering
\begin{tikzpicture}[x=0.69cm,y=0.58cm,>=Stealth]
  \draw[fill=black!8] (1,2) rectangle (6,2.5);
  \draw[fill=black!8] (2,1) rectangle (7,1.5);
  \draw[fill=black!8] (3,0) rectangle (8,0.5);
  \node[anchor=east] at (0.75,2.25) {seed 1};
  \node[anchor=east] at (0.75,1.25) {seed 2};
  \node[anchor=east] at (0.75,0.25) {seed 3};
  \draw[densely dashed] (1,-0.2) -- (1,-1.5);
  \draw[densely dashed] (8,-0.2) -- (8,-1.5);
  \draw[very thick] (1,-1) -- (11,-1);
  \draw[->] (11.25,-1) -- (12,-1);
  \node[above] at (4.5,-1) {seed union agrees with $\pi$};
  \node[below] at (9.5,-1) {continuation};
  \draw (10.88,-1.13) -- (11.12,-0.87);
  \draw (10.88,-0.87) -- (11.12,-1.13);
  \node[above] at (11,-0.75) {$q$};
  \node[below] at (11.1,-1.55) {first mismatch};
\end{tikzpicture}
\caption{Different seed ends initiate the same first-mismatch test.
The overlapping seed union already agrees with the common repetition,
so no scan can stop inside a later seed.}
\label{fig:common-seeds}
\end{figure}

For an outer layer with full cut-offset set $\mathcal H_i$, a suffix
beginning at $k$ uses
\[
 \mathcal T=\{k-h:h\in\mathcal H_i\}.
\]
If the prefix chooses an actual cut $k=t+h_0$ with
$h_0\in\mathcal H_i$, the true origin belongs to this set regardless of
the suffix's independent role or tag.  Its diameter is
$\operatorname{diam}\mathcal H_i$.  Thus the end guard in
Lemma~\ref{lem:shared-anchor} establishes exact consumption before any
agreement of tags is used.  Likewise an aligned stencil $\mathcal P_i$
gives candidate origins $k+L_0-p'$; at a true cut $k=s+p$ these are
$t+p-p'$.  Their diameter is $\operatorname{diam}\mathcal P_i$.

\begin{lemma}[End selection after bounded decoding]\label{lem:decoded-end}
Let $e,k,[k,b),t,q$ be as in Lemma~\ref{lem:shared-anchor}, and let
$\mathcal T$ be any finite candidate set containing $t$, with no diameter
restriction.  Suppose its candidates are tested by bounded predicates
which do not use an anchor, all their inspection intervals lie inside
both the actual episode and the proposed suffix, and the true origin
passes its predicate.  If exactly one candidate passes, it is the true
origin.  Applying the exact end rule at that decoded origin, with its
search and seed data also present in both factors, forces the proposed
suffix to end at $q+B_{\rm tail}$ without any stability assumption.
\end{lemma}

\begin{proof}
The bounded predicates have the same values on their common letters,
so the passing true candidate cannot disappear in the proposed suffix.
Uniqueness therefore identifies it.  Apply
Lemma~\ref{lem:shared-anchor} with the singleton reference set $\{t\}$;
its equality condition is automatic.
\end{proof}

The last aligned stage will use Lemma~\ref{lem:decoded-end}.  Its current
witness is a bounded end-marker test independent of $q$, so decoding
precedes the end rule.  At the other stages one reference end rule is
applied first.  In either order, subsequent controls at translated
origins keep that same $q$ fixed.  They do not apply additional
first-mismatch searches.

\subsection{The finite margin interface}

The end guards require every bounded inspection to lie inside the true
episode, even when the suffix proposes a wrong origin or a wrong end.
This requirement can be met after all finite schedules have been chosen.
We state the margin calculation here; Proposition~\ref{prop:parameters}
will supply the complete inventory of inspections for the actual splits.

Take the union of the bounded inspection intervals under every allowed
cut, metadata choice, and candidate origin. Include the neighborhoods
needed to decide markers and the intervals swept out in
Lemma~\ref{lem:search-crossings}. Let $I_-\le0\le I_+$ bound the
inner-marker and trial data and all outer cuts relative to $t$, and let $E_-,E_+$ bound
the end-marker inspections and seed blocks relative to $z_\bullet(t)$.
These are finite bounds. The continuation to the selected first mismatch
is governed separately by the exact end rule; products ending there and
residues involving $q$ do not enlarge this bounded inventory.

Let $A$ bound the offsets from $s$ of the buffers before $t$.
Sufficient final margin inequalities are
\begin{equation}\label{eq:geometry-margins}
\begin{split}
 L_0+I_-&>A,\\
 o_\bullet+\min\{E_-,-d_\bullet\}&>I_+,\\
 B_{\rm tail}&>E_++d_\bullet,
 \qquad B_{\rm tail}\ge1.
\end{split}
\end{equation}
The first inequality puts all inner data after the early buffers.
The second puts the end controls and every possible anchor after the
inner data and the outer cuts included in that inventory. Since
$q\ge z_\bullet(t)-d_\bullet$, the last inequality puts every end
control strictly before even the earliest complete end
$z_\bullet(t)-d_\bullet+B_{\rm tail}$. Thus all bounded controls
are internal to every true complete episode.

Internal visibility does not by itself put data on the correct side of
a cut. Each split must also place its prefix computations before its cut
and its suffix computations after it. Any supplied left context must be
declared as finite metadata and verified by the prefix. The schedules in
Proposition~\ref{prop:parameters} will meet these reading-side conditions.

Every test argument must itself contain every position it is required
to inspect on its reading side. An unavailable candidate or translate
causes rejection; it is never omitted from a universal guard or uniqueness
test. Requested product intervals must likewise have available, correctly
ordered endpoints. This applies to a standalone suffix with false metadata
and to a proposed endpoint different from the true episode end. The margin
inequalities give availability for the true suffix; the rejection rule
ensures that every accepted suffix sees the same bounded control values
used in Lemmas~\ref{lem:shared-anchor} and~\ref{lem:decoded-end}.

The remaining obligation is to fix this finite inventory independently of
$L_0,o_\bullet,B_{\rm tail}$. Proposition~\ref{prop:parameters} does so
before choosing those three lengths in the order displayed above.
